\begin{mistake}{2026-04-12}{22}

\begin{problemcontent}
计算极限：\(\displaystyle \lim_{x \to 0} \left[ \frac{1+\int_0^x \sin t \,\mathrm{d}t}{\ln(1+x^2)} - \frac{1}{\sin^2 x} \right]\)
\end{problemcontent}

\begin{solutioncontent}
先利用牛顿-莱布尼茨公式计算变上限积分：
\(\int_0^x \sin t \,\mathrm{d}t = [-\cos t]_0^x = 1 - \cos x\)。
代入原极限并通分：
\[
\begin{aligned}
L &= \lim_{x \to 0} \left[ \frac{2 - \cos x}{\ln(1+x^2)} - \frac{1}{\sin^2 x} \right] \\
&= \lim_{x \to 0} \frac{(2-\cos x)\sin^2 x - \ln(1+x^2)}{\ln(1+x^2) \sin^2 x}
\end{aligned}
\]

对于分母，利用等价无穷小 \(\ln(1+x^2) \sim x^2\)，\(\sin^2 x \sim x^2\)，故分母 \(\sim x^4\)。
这说明分子需用泰勒展开至 \(x^4\) 阶：
\[
\cos x = 1 - \frac{x^2}{2} + \frac{x^4}{24} + o(x^4) \implies 2-\cos x = 1 + \frac{x^2}{2} - \frac{x^4}{24} + o(x^4)
\]
\[
\sin^2 x = \left(x - \frac{x^3}{6} + o(x^3)\right)^2 = x^2 - \frac{x^4}{3} + o(x^4)
\]
\[
\ln(1+x^2) = x^2 - \frac{x^4}{2} + o(x^4)
\]

将展开式代入分子：
\[
\begin{aligned}
\text{分子} &= \left( 1 + \frac{x^2}{2} \right)\left( x^2 - \frac{x^4}{3} \right) - \left( x^2 - \frac{x^4}{2} \right) + o(x^4) \\
&= \left( x^2 - \frac{x^4}{3} + \frac{x^4}{2} \right) - x^2 + \frac{x^4}{2} + o(x^4) \\
&= \left(\frac{1}{2} - \frac{1}{3} + \frac{1}{2}\right)x^4 + o(x^4) = \frac{2}{3}x^4 + o(x^4)
\end{aligned}
\]

因此，原极限：
\[
L = \lim_{x \to 0} \frac{\frac{2}{3}x^4}{x^4} = \frac{2}{3}
\]
\end{solutioncontent}

\mistakeknowledge{
\begin{itemize}
  \item \textbf{核心公式}：麦克劳林展开式：\(\cos x = 1 - \frac{x^2}{2!} + \frac{x^4}{4!} + o(x^4)\)，\(\ln(1+x) = x - \frac{x^2}{2} + o(x^2)\)。
  \item \textbf{易错点}：加减法中不可直接使用等价无穷小替换；泰勒展开的阶数必须一致对齐到 \(x^4\)，容易漏掉 \(\sin^2 x\) 中 \(\frac{x^4}{3}\) 这一项。
  \item \textbf{关联知识}：看到通分后分母为 \(x^4\) 量级，分子的所有因式乘积和加减都必须精确保留到 \(4\) 阶导数级别。
\end{itemize}}

\end{mistake}
